% Part: first-order-logic
% Chapter: sequent-calculus
% Section: derivations

\documentclass[../../../include/open-logic-section]{subfiles}

\begin{document}

\iftag{FOL}
      {\olfileid{fol}{seq}{der}}
      {\olfileid{pl}{seq}{der}}

\olsection{\usetoken{P}{derivation}}

\begin{explain}
اسیں دَس چکے آں کہ ابتدائی سیکوئنٹ کیہو جہا ہُندا اے، تے استنباط دے
قاعدے وی دے چکے آں۔ سیکوئنٹ حساب وچ !!^{derivation}s ایہناں توں
استقرائی طور تے بنائے جاندے نیں: ہر !!{derivation} یا تاں اپنے آپ وچ
اک ابتدائی سیکوئنٹ ہُندا اے، یا اک یا دو !!{derivation}s تے اوہناں دے
بعد اک استنباط اُتے مشتمل ہُندا اے۔
\end{explain}

\begin{defn}[$\Log{LK}$ !!{derivation}]
اک \emph{$\Log{LK}$-!!{derivation}} جیہڑی سیکوئنٹ~$S$ دی اے،
سیکوئنٹاں دا اک محدود درخت اے جیہڑا ہِٹھ دتیاں شرطاں پوری کردا اے:
\begin{enumerate}
\item درخت دے سب توں اُتلے سیکوئنٹ ابتدائی سیکوئنٹ ہُندے نیں۔
\item درخت دا سب توں ہِٹھلا سیکوئنٹ~$S$ اے۔
\item درخت وچ $S$ توں سوا ہر سیکوئنٹ، استنباط دے قاعدے دی کسے درست
  ورتوں دا مقدمہ اے، جیہدا نتیجہ درخت وچ اوس سیکوئنٹ دے عین ہِٹھاں
  کھلوتا اے۔
\end{enumerate}
فیر اسیں $S$ نوں !!{derivation} دا \emph{آخرلا سیکوئنٹ} آکھدے آں تے
ایہ وی آکھدے آں کہ $S$، \emph{$\Log{LK}$ وچ !!{derivable} اے} (یا
$\Log{LK}$-!!{derivable} اے)۔
\end{defn}

\begin{ex}
ہر ابتدائی سیکوئنٹ، مثلاً $!C \Sequent !C$، اک !!a{derivation} اے۔
اس اُتے، من لو، $\LeftR{\Weakening}$ قاعدہ لا کے اسیں اک نویں
!!{derivation} حاصل کر سکدے آں،
\begin{prooftree}
\Axiom$ \Gamma \fCenter \Delta $
\RightLabel{\LeftR{\Weakening}}
\UnaryInf$ !A, \Gamma \fCenter \Delta$
\end{prooftree}
پر ایہ قاعدہ عام روپ وچ ورتن لئی اے: قاعدے وچلے $!A$ دی تھاں اسیں
کوئی وی !!{sentence} رکھ سکدے آں، مثلاً~$!D$ وی۔ جے مقدمہ ساڈے
ابتدائی سیکوئنٹ $!C \Sequent !C$ نال میل کھاندا اے، تاں ایس دا مطلب
اے کہ $\Gamma$ تے $\Delta$ دونیں صرف~$!C$ نیں، تے نتیجہ فیر
$!D, !C \Sequent !C$ ہووے گا۔ سو ہِٹھ لکھی شے اک !!a{derivation} اے:
\begin{prooftree}
\Axiom$ !C \fCenter !C $
\RightLabel{\LeftR{\Weakening}}
\UnaryInf$ !D, !C \fCenter !C$
\end{prooftree}
ہُن اسیں کوئی ہور قاعدہ، من لو $\LeftR{\Exchange}$، لا سکدے آں،
جیہڑا سانوں کھبے پاسے دے دو !!{sentence}s نوں آپو وچ بدل دین دی
اجازت دیندا اے۔ سو ہِٹھ لکھی شے وی اک درست !!{derivation} اے:
\begin{prooftree}
\Axiom$ !C \fCenter !C $
\RightLabel{\LeftR{\Weakening}}
\UnaryInf$ !D, !C \fCenter !C$
\RightLabel{\LeftR{\Exchange}}
\UnaryInf$ !C, !D \fCenter !C$
\end{prooftree}
ایس قاعدے دی ورتوں وچ—قاعدہ ایویں دِتا گیا سی—
\begin{prooftree}
\Axiom$ \Gamma, !A, !B, \Pi \fCenter \Delta $
\RightLabel{\LeftR{\Exchange}}
\UnaryInf$ \Gamma, !B, !A, \Pi \fCenter \Delta,$
\end{prooftree}
$\Gamma$ تے $\Pi$ دونیں خالی سن، $\Delta$، $!C$ اے، تے $!A$ تے
$!B$ دے کردار ترتیب وار $!D$ تے~$!C$ نے نبھائے نیں۔ لگ بھگ
ایسے طریقے نال اسیں ایہ وی ویکھدے آں کہ
\begin{prooftree}
\Axiom$ !D \fCenter !D $
\RightLabel{\LeftR{\Weakening}}
\UnaryInf$ !C, !D \fCenter !D$
\end{prooftree}
اک !!a{derivation} اے۔ ہُن اسیں ایہناں دوناں !!{derivation}s نوں لے
کے $\RightR{\land}$ نال اکٹھیاں کر سکدے آں۔ اوہ قاعدہ ایہ سی
\begin{prooftree}
\Axiom$\Gamma \fCenter \Delta, !A$
\Axiom$ \Gamma \fCenter \Delta, !B$
\RightLabel{\RightR{\land}}
\BinaryInf$ \Gamma \fCenter \Delta, !A \land !B $
\end{prooftree}
ساڈے معاملے وچ، مقدمیاں نوں اوہناں !!{derivation}s دے آخرلے سیکوئنٹاں
نال میل کھانا چاہیدا اے جیہڑے ایہناں مقدمیاں اُتے ختم ہُندے نیں۔ ایس دا مطلب
اے کہ $\Gamma$، $!C, !D$ اے، $\Delta$ خالی اے، $!A$، $!C$ اے تے
$!B$، $!D$ اے۔ سو جے استنباط درست ہونا اے، تاں نتیجہ $!C, !D \Sequent !C
\land !D$ اے۔
\begin{prooftree}
\Axiom$ !C \fCenter !C $
\RightLabel{\LeftR{\Weakening}}
\UnaryInf$ !D, !C \fCenter !C$
\RightLabel{\LeftR{\Exchange}}
\UnaryInf$ !C, !D \fCenter !C$
\Axiom$ !D \fCenter !D $
\RightLabel{\LeftR{\Weakening}}
\UnaryInf$ !C, !D \fCenter !D$
\RightLabel{\RightR{\land}}
\BinaryInf$ !C, !D \fCenter !C \land !D $
\end{prooftree}
یقیناً، اسیں مقدمیاں دی ترتیب اُلٹ وی سکدے آں؛ فیر $!A$، $!D$
ہووے گا تے $!B$،~$!C$ ہووے گا۔
\begin{prooftree}
\Axiom$ !D \fCenter !D $
\RightLabel{\LeftR{\Weakening}}
\UnaryInf$ !C, !D \fCenter !D$
\Axiom$ !C \fCenter !C $
\RightLabel{\LeftR{\Weakening}}
\UnaryInf$ !D, !C \fCenter !C$
\RightLabel{\LeftR{\Exchange}}
\UnaryInf$ !C, !D \fCenter !C$
\RightLabel{\RightR{\land}}
\BinaryInf$ !C, !D \fCenter !D \land !C $
\end{prooftree}
\end{ex}

\end{document}
